% !TeX root = ../papers/nc-p3.tex
\section{薄い係数が非零で交差ブロックの階数が一の両対称型}
\label{sec:klein-singular-exchange-corners}
\begin{theorem}\label{thm:klein-singular-exchange-corners}
定理\ref{thm:klein-bisymmetric-irreducible}の固定色で $a_ib_j\ne0$ が全ての $i,j$ に対して成り立ち、
$\det U=0$ とする。正則なものは全て、文字の付け替えを許して
\[
 M_0=\begin{pmatrix}X&X&1&1\\Y&Y&1&1\\1&1&W&W\\1&1&V&V\end{pmatrix},
 \qquad M_1=M_0^{\mathsf t},\qquad M_2=M_3=\boldsymbol1_{4\times4},
 \qquad (X-Y)(W-V)\ne0
\]
で与えられる。この全域で正表示と非自明捻りの双方が標準Koszul左右AS正則で、最小周期は1である。
\end{theorem}
\begin{proof}
偶奇の文字倍率を $e_i,f_i$ とすると、薄い因子 $a_i,b_i$ への作用は
$e_if_{i+2},e_if_{i+3}$ である（文字添字の加法は四元群）。因子の共通逆倍率を選んで
$a_0a_1=b_0b_1$ とでき、$\det U=0$ より $a_2a_3=b_2b_3$ も成り立つ。
すると各組で $f_{i+2}/f_{i+3}=b_i/a_i$ を解いて $a_i=b_i=1$ にできる。
$J=\boldsymbol1_{2\times2}$ とおいて前の階数分解で $LD'=RC=J$ を得る。

$L,C$ が共に階数1なら、$L$ の列と $C$ の行は定数ベクトルの倍数になる。
関係の左評価を $x=(1,1,0,0)$ で取れば、厚い二色は低位側の $(1,1)$ の方向だけを、
薄い二色は高位側の $(1,1)$ の方向だけを与える。定理\ref{thm:all-klein-bisymmetric-evaluation}
に反する。$R,D'$ が共に階数1の場合は高位二文字で同様である。
$L,D'$ の両方または $R,C$ の両方が可逆であることは $J$ の階数から除かれる。
従って $L,R$ が共に可逆、または $C,D'$ が共に可逆である。後者は全ての奇文字を1だけ平行移動し、
$M_0$ と $M_1=M_0^{\mathsf t}$ を交換して前者に帰着する。
前者では $C,D'$ の各二列が一致し、$L,R$ の可逆性が二つの差の非消滅に等しい。

元の総文字1の表示で第一辺の関係空間を $S$ とする。対角色では
$(X,Y,1,1),(1,1,W,V)$、色1では対称積 $01,23$、残る二色では対称積の和
$02+13,03+12$ が基底となる。ここでは異なる二文字の積を対称化して記した。
両対称性から隣の関係は $S^{\mathsf t}=S$ であり、全辺が同じになる。
$tt^{\mathsf t}\in S$ で両半分が非零なら、交差ブロックは対角成分同士、
非対角成分同士が等しい階数1の行列だから $t_0^2=t_1^2,t_2^2=t_3^2$。
対角の二係数がともに零となり矛盾する。片方が零の場合も、対角二平面から各半分への射影が
同型なので非零解はない。従って可換二次双対の四式は共通射影零点を持たず、正則列である。
対称チャートの定理から全次数の結論を得る。
非自明捻りには同じ四文字次数で全辺の関係を捻ればよい。実際テンソルの四位置を平行移動して
総文字を零にした時の符号 $\prod_{i<j}\mu(q_i,q_j)$ の変化は、双指標性により
各位置の文字倍率の積である。従って同じ斉次代数のコサイクル捻りがindexed表示を与え、
周期1も保たれる。
\end{proof}

\begin{theorem}\label{thm:singular-exchange-negative-quotient}
上の捻りの印を忘れた指定始点の商は $\mathbb A^2\times\mathbb G_m$ である。座標は
\[
 u=(X+Y)/(X-Y),\quad v=(W+V)/(W-V),\quad h=(X-Y)(W-V),
 \qquad (s,\ell,H)=(u^2+v^2,uvh,h^2).
\]
反対階数は $16-2[u=0]-2[v=0]$、射影安定群の位数は通常4、$u=\pm v$ で8、
$u=v=0$ で16である。この三次元族の閉包は四次元負成分の周期1点の別の因子を与える。
\end{theorem}
\begin{proof}
$\mu(i,j)=(-1)^{i_{\rm hi}j_{\rm lo}}$ とする。捻りの点対 $(x,y)$ で $x$ の両半分が非零なら、
色1の式より $y=(m x_0,-m x_1,n x_2,n x_3)$。薄い二式は
$(m+n)(x_0x_2+x_1x_3)=(m+n)(x_0x_3+x_1x_2)=0$。
後ろの二因子が両方零なら $x_1=\pm x_0,x_3=\mp x_2$ であり、残りの二式から
$m(X-Y)=n(W-V)=0$ となる。従って $n=-m$ で、点を与える条件は
\[
 Xx_0^2-Yx_1^2=x_2^2-x_3^2,\qquad x_0^2-x_1^2=Wx_2^2-Vx_3^2 .
\]
低位・高位それぞれで二列が独立だから、鉛筆の四列の比例は異なる半分に属する列の間にしか起きない。
従って滑らかな四次曲線、異なる二平面二次曲線、または四直線の四角形を得る。
平方比例式を因数分解すれば後二者も被約である。片半分が零なら色1の式と対角二式で
同じ半分の $y$ は零になり、残る二式から孤立点 $(1,\pm1,0,0),(0,0,1,\pm1)$ を得る。
右側でも両半分の非零な部分は同じ符号変更のグラフで、孤立点は同じ二対である。
定理\ref{thm:negative-pairing-unmarked-quotient}の四頂点の復元が特異四次でも適用される。
すなわち二次曲線二本なら鉛筆の二単頂点を含む孤立点の割線を使い、
四角形なら節点の二対角線の両方と交わる孤立点の割線を使う。

四頂点の枠で総次数零に戻すと、薄い48語の係数は $\pm1$ である。同定できる置換は
前と同じ32候補であり、薄い係数を合わせた後の対角残余は
$(X,Y,W,V)\mapsto(kX,kY,k^{-1}W,k^{-1}V)$ である。
実際、厚い各半分内の二語の文字指数の差と異なる半分から取った二語の指数の和が
薄い48語の整数行格子に入る。逆に上位ビットの個数 $b$ による倍率 $k^{1-b/2}$ は実現する。
四位置の置換補正を薄い文字指数行列のSmith形で解き、残り十六語を比較すると、
八候補は両組内の独立な交換と二組の交換を与え、他の二十四候補には
$X=Y,W=V$ が必要である。この計算は厚い係数で割らずに行える。
一つの低位語を足した指数行列のSmith因子は $1^{11},2,2$ なので、
各実現置換につき射影対角安定群の位数は4である。
\path{computations/verify_exchange_singular_corners.py} は整数行格子、全置換補正および反対切断を検算する。
反対切断の二色は常に可逆で、他の各色は $X+Y,W+V$ の消滅ごとに階数を一つ失う。

共通倍率を除くと $(u,v,h)$ となり、各組内の交換は対応する $u$ または $v$ と $h$ の
符号を同時に変える。これら二交換の不変環は $u^2,v^2,uvh,h^{\pm2}$ で生成され、
$(uvh)^2=u^2v^2h^2$。二組の交換は $u,v$ を交換するので主張の商となる。
安定群の位数もこの作用から従う。負成分への包含は定理\ref{thm:klein-bisymmetric-irreducible}による。
以前の $zD(z^2-1)\ne0$ の表示のうち四孤立点を持つ部分とは、薄い交差比がここでは1、
先の薄い置換では $z^{\pm1}$ であるため一致しない。他の部分には四孤立点がない。
従ってこの既約三次元像は先の周期1因子の稠密開集合と交わらず、閉包も異なる。
\end{proof}

\begin{proposition}\label{prop:singular-corner-positive-forms}
正表示の印を忘れた同型には次の有限な正規形集合が使える。対称関係の直交空間を四次元の
対称行列空間 $L$ とし、階数2の二元を射影的に選んで各々の二つの線型因子を取る。
四因子が独立な全ての組を全ての順序で基底として、$L$ のPlücker座標の台と対角トーラスの
文字不変量を記録する。この有限集合が交わることが同型の必要十分条件である。
\end{proposition}
\begin{proof}
$L$ の行列は
\[
 \begin{pmatrix}A&0&r&t\\0&B&-t&-r\\r&-t&C&0\\t&-r&0&D\end{pmatrix},\qquad
 XA+YB+C+D=A+B+WC+VD=0 .
\]
$r=t=0$ の二次元空間で $(A,B)$ と $(C,D)$ はいずれも座標になる。
三次小行列には符号を除き $rt(A+B),rt(C+D)$ がある。$rt\ne0$ なら
階数2以下では全対角が零、$r=\pm t$ である。この二元はそれぞれ低位と高位の
二文字の和または差を因子とし、合わせて四因子は独立である。
$t=0,r\ne0$ なら $AC=BD=r^2$、$r=0,t\ne0$ なら $AD=BC=t^2$ である。
対角直線上で $AC-BD$ は恒等的に零ではない。恒等的なら独立な一次式 $A,B$
から $C=kB,D=kA$ が従い、第一式より $X=Y$ となる。$AD-BC$ も同様である。
$r=t=0$ の階数2点もどちらかの式に含まれる。従って各場合高々四点、先の二点と合わせて高々十点である。
階数1の元は、非零の非対角成分がある場合は二つのブロックの階数から、対角だけなら
両射影の独立性から除かれる。各枠での文字検査は、台上のPlücker単項式の重み（全体倍率の列も含む）
の整数関係の基底を取り、対応する単項式比を記録すればよい。系\ref{cor:symmetric-chart-isomorphism}
により同型は正確に $L$ の合同であり、合同は因子枠を運ぶので判定は必要十分となる。
同じ検査は任意の他の対称表示との比較にも使え、階数2点の数または独立枠の存在が異なれば直ちに分離する。
\end{proof}
